\scanpage{137}
\mspage{53}

de deux foncteurs adjoints $(f^{*},f_{*})$ et des morphismes d'adjonction. Quand $f^{*}$ est une équivalence, ce qui équivaut à $f_{*}$ une équivalence, ou $\alpha$ et $\beta$ des iso, on dira qu'on a une \uline{équivalence de pseudo-topos}, et à toutes fins pratiques, on pourra alors regarder $\unsure{X|\mathcal X}$ et $\unsure{Y|\mathcal Y}$ comme des objets essentiellement identiques.

\deleted{\illegible{1}} Pour deux pseudo-topos $X,Y$, j'ai envie d'introduire (comme dans le cas des topos) un pseudo-topos produit

\[
X\times Y,
\tag*{(4.1.4)}
\]

% slanted
\margin{Il apparaît finalement que ce qu'on construit n'est \uline{pas} un 2-produit \unsure{! (dans la|(dans la} 2-catégorie envisagée.) \deleted{\unsure{donc}} il vaut mieux le noter par un signe tel que $X\boxtimes Y$. Mais si $X,Y$ sont deux topos, $X\boxtimes Y$ est un 2-produit dans la 2-catégorie des topos.}

ayant les propriétés d'un 2-produit dans la 2-catégorie des pseudo-topos. \unsure{mais|Mais} précisons que par définition

\[
\begin{aligned}
\underline{\operatorname{Hom}}_{\mathrm{ps\,top}}(X,Y)
&\simeq\underline{\operatorname{Hom}}_{!}(F(Y),F(X))^{\mathrm{o}},\\
&\simeq\underline{\operatorname{Hom}}^{!'}(F(X),F(Y)),
\end{aligned}
\tag*{(4.1.5)}
\]

où l'\deleted{exposant}\inserted{indice} ! \inserted{(resp. l'exposant !)} indique la sous-catégorie pleine de la catégorie $\underline{\operatorname{Hom}}$ envisagée, formée des foncteurs qui commutent aux \inserted{petites} $\underset{\longrightarrow}{\lim}$ (resp. aux petites
% archive sheet number: 136
